% !TeX program = lualatex
% Definition and sources only; this file is not an LLM solution attempt.
\documentclass[11pt,a4paper]{article}
\usepackage[margin=25mm]{geometry}
\usepackage{fontspec}
\setmainfont{lmroman10-regular.otf}[BoldFont=lmroman10-bold.otf,
 ItalicFont=lmroman10-italic.otf,BoldItalicFont=lmroman10-bolditalic.otf]
\usepackage{amsmath,amssymb,mathtools}
\usepackage{unicode-math}
\setmathfont{latinmodern-math.otf}
\AtBeginDocument{\let\setminus\smallsetminus}
\usepackage{xurl,hyperref}
\hypersetup{colorlinks=true,urlcolor=blue}
\setcounter{secnumdepth}{0}
\setlength{\parindent}{0pt}
\setlength{\parskip}{5pt}
\setlength{\emergencystretch}{3em}
\begin{document}
\section*{Research Problems in Function Theory — Problem 2.88}\label{statement}

\subsection{Definitions and mathematical statement}
For $c\in\mathbb C$, write $P_c(z)=z^2+c$. The Mandelbrot set is
\[
\mathcal M=\{c\in\mathbb C:\sup_{n\ge0}|P_c^n(0)|<\infty\},
\]
and let $B=\partial\mathcal M$ be its topological boundary in the parameter plane. These are sets of parameters $c$, rather than Julia sets in the dynamical $z$-plane.
The remaining questions in this source item are:
\begin{enumerate}
\item Is $B$ locally connected? Precisely, for every $b\in B$ and every neighbourhood $U$ of $b$ in the relative topology of $B$, must there be a connected relative neighbourhood $V$ with $b\in V\subset U$?
\item Does $B$ have two-dimensional Lebesgue measure zero?
\end{enumerate}
The original item also asked whether the Hausdorff dimension of $B$ is two. The source update records Shishikura's proof of that assertion, so it is supplied as established context rather than a further requested conjecture. Full Hausdorff dimension alone does not decide whether planar area vanishes. Local connectedness is asked at every boundary point, including infinitely renormalizable parameters; results for special parameter classes do not settle this unrestricted formulation.
\subsection{Short English statement}
Is the Mandelbrot-set boundary locally connected everywhere, and does it have zero planar area?
\subsection{Sources}
\begin{itemize}
\item[S1] ULAM.AI and the UnsolvedMath Contributors, supplied problems.json, catalogue entry, AMR Open Problem Lists. Catalogue identity and source transcription; CC BY 4.0.
\url{https://huggingface.co/datasets/ulamai/UnsolvedMath}.
\item[S2] Hayman - Research Problems in Function Theory (2018), Problem 2.88. 
\url{https://arxiv.org/abs/1809.07200}.
\end{itemize}
\end{document}
